\usepackage[utf8]{inputenc}
\usepackage[T1]{fontenc}
\usepackage{csquotes}
\usepackage{graphicx}
\usepackage{amsmath}
\usepackage{booktabs}
\usepackage{hyperref}
\usepackage{acronym} % geladen von twbook.cls; nohyperlinks via \PassOptionsToPackage in mt_bhinder.tex
\usepackage[most]{tcolorbox}
\usepackage{listings}
\usepackage{float} % erlaubt [H] fuer harte Float-Platzierung
\usepackage{longtable} % seitenumbrechende Tabellen
\usepackage{pdflscape} % echte Querformat-Seiten (PDF-Viewer zeigt horizontal)
% Absatzstil: kein Erstzeilen-Einzug, stattdessen vertikaler Abstand
% zwischen Absaetzen (Blockstil wie in der Salma-Arbeit).
\usepackage{parskip}

% Bibliografie: kein haengender Einzug in den Folgezeilen, dafuer
% kleiner Abstand zwischen den Eintraegen.
\setlength{\bibhang}{0pt}
\setlength{\bibitemsep}{0.6\baselineskip}

% Sentence Case auch fuer die generierten Verzeichnis-Titel
% (Konsistenz mit den Kapitel-/Section-Ueberschriften).
\addto\captionsenglish{%
    \renewcommand*{\listfigurename}{List of figures}%
    \renewcommand*{\listtablename}{List of tables}%
}
\renewcommand*{\listacroname}{List of abbreviations}
\graphicspath{{images/}}

% Stil fuer Prompt-/Code-Listings (z.B. lst:system-prompt in Kap. 3.5)
\lstset{
    basicstyle=\ttfamily\small,
    breaklines=true,
    frame=single,
    framerule=0.5pt,
    xleftmargin=2mm,
    captionpos=b}

% Box fuer Forschungsfragen (Kap. 1) und deren Beantwortung (Kap. 4).
% Schlichter Rahmen: duenne graue Linie, weisser Koerper, hellgraue
% Titelzeile, eckige Kanten.
% Box fuer Antwort-Auszuege in Kap. 4 (gleiche Optik wie rqbox).
\newtcolorbox{responsebox}[1]{%
    breakable,
    sharp corners,
    colback=gray!2,
    colframe=black!60,
    colbacktitle=black!8,
    coltitle=black,
    fonttitle=\bfseries\small,
    fontupper=\small,
    title={#1},
    boxrule=0.5pt,
    left=2mm, right=2mm, top=1.5mm, bottom=1.5mm}

\newtcolorbox{rqbox}[1]{%
    breakable,
    sharp corners,
    colback=white,
    colframe=black!60,
    colbacktitle=black!8,
    coltitle=black,
    fonttitle=\bfseries,
    title={#1},
    boxrule=0.5pt,
    left=2mm, right=2mm, top=1.5mm, bottom=1.5mm}